\begin{corollary}[6.14.5]
\label{IV.2.s6.14.5-fr}
Sean $A$ un anillo noetheriano y $A'$ una $A$-álgebra noetheriana
tal que el morfismo $\operatorname{Spec}(A')\to\operatorname{Spec}(A)$
sea normal. Sean $B$ y $C$ dos $A$-álgebras,
$\varphi:B\to C$ un $A$-homomorfismo que hace de $C$ una
$B$-álgebra, y $D$ la clausura integral de $B$ en $C$. Si se pone
$B'=B\otimes_A A'$, $C'=C\otimes_A A'$ y $D'=D\otimes_A A'$,
entonces $D'$ es la clausura integral de $B'$ en $C'$.
\end{corollary}

\begin{proof}
Sea $(B_\lambda)$ la familia filtrante creciente de las
sub-$A$-álgebras de tipo finito de $B$, y pongamos
$B'_\lambda=B'_\lambda\otimes_A A'$. Así, $B'$ es la unión de la
familia filtrante creciente de las sub-$A'$-álgebras de tipo finito
$B'_\lambda$. Sea $D_\lambda$ la clausura integral de $B_\lambda$
en $C$, y pongamos $D'_\lambda=D_\lambda\otimes_A A'$; entonces
$D$ es la unión de los $D_\lambda$ y $D'$ la de los $D'_\lambda$.
Si demostramos que $D'_\lambda$ es la clausura integral de
$B'_\lambda$ en $C'$, se seguirá que $D'$ es la clausura integral
de $B'$ en $C'$. Ahora bien, $B_\lambda$ y $B'_\lambda$ son
noetherianos, y el morfismo
$\operatorname{Spec}(B'_\lambda)\to\operatorname{Spec}(B_\lambda)$
es normal (6.8.2). Por tanto, se puede aplicar (6.14.4) sustituyendo
$A$, $A'$ y $B$ por $B_\lambda$, $B'_\lambda$ y $C$,
respectivamente, de donde resulta el corolario.
\end{proof}

Por ejemplo, se puede aplicar (6.14.5) cuando $A$ es un anillo \emph{local excelente} y $A'$ es su completado $\widehat A$, puesto que
en este caso $\operatorname{Spec}(A')\to\operatorname{Spec}(A)$
es un morfismo regular (7.8.2).
